\section{Definitions and the foundational inequality}

All graphs are finite and nonempty. An oriented graph has no loops and
no pair of oppositely directed arcs. Write $\Nout(x)$ and $\Nin(x)$ for
its out- and inneighborhoods. The sets $\Ntwo(x)$ and $\Nback(x)$ consist,
respectively, of vertices at directed distance exactly two from $x$ and
to $x$. In particular, direct neighbors are excluded. Their cardinalities
are denoted by $\dout(x),\din(x),\dtwo(x),\dback(x)$.

A Seymour vertex satisfies $\dtwo(x)\ge\dout(x)$. Call a graph
\emph{all-deficient} if it has no Seymour vertex, and set
\[
g(x)=\dout(x)-\dtwo(x)>0.
\]
An all-deficient graph is \emph{arc-minimal} here if deleting any nonempty
set of arcs, with the vertex set fixed, destroys this property.
This is stronger than merely assuming that each single arc is
individually indispensable. Every all-deficient graph has an
arc-minimal spanning subgraph.

If a counterexample exists, take an arc-minimal spanning subgraph and
then a sink strongly connected component. That component is still
all-deficient and inherits arc-minimality: an arc deletion confined to it
cannot improve any vertex outside it. Thus the strongly connected class
considered here is nonempty whenever the conjecture is false.

Let $\mathcal M(x)$ be the set of vertices distinct from and
nonadjacent to $x$, let $\mu(x)=|\mathcal M(x)|$, and let $M$ be the
total number of missing unordered pairs. Thus $\sum_x\mu(x)=2M$.

\begin{theorem}\label{thm:main}
Let $D$ be a strongly connected arc-minimal all-deficient oriented graph
on $n$ vertices. Then, for every vertex $x$,
\begin{equation}\label{eq:main}
\dtwo(x)-\dback(x)\le\mu(x)-1.
\end{equation}
Consequently,
\begin{equation}\label{eq:global}
M\ge\left\lceil\frac n2\right\rceil,\qquad
|A(D)|\le\left\lfloor\frac{n(n-2)}2\right\rfloor.
\end{equation}
\end{theorem}

The restriction to arc-minimal graphs matters. Deleting arcs increases
the number of missing pairs, so \eqref{eq:global} cannot be transferred
backwards to an arbitrary denser counterexample. We do not claim that
every oriented graph with fewer than $n/2$ missing pairs satisfies the
conjecture.


\section{An external-boundary deletion lemma}

For $S\subseteq V(D)$, put
\[
\Gamma(S)=\bigcup_{v\in S}\Nout(v),\qquad
\bd S=\Gamma(S)\setminus S,\qquad
\bd^2S=\Gamma(\bd S)\setminus(S\cup\bd S).
\]
These are external neighborhoods. Seacrest's notation~\cite{SeacrestSubsets2019} permits
intersection with the starting set, so the following adapted statement
and proof are made explicit. The proof is included to make the external convention explicit.

\begin{lemma}\label{lem:boundary}
In an arc-minimal all-deficient graph, every $S$ with $\bd S\ne\varnothing$
satisfies $|\bd^2S|<|\bd S|$.
\end{lemma}
\begin{proof}
Choose an inclusion-maximal nonempty $T\subseteq\bd S$ with
$|\bd T\setminus S|<|T|$, or $T=\varnothing$ if none exists.
If $T=\bd S$, the assertion follows. Otherwise delete all arcs
$S\to\bd S\setminus T$. Any resulting Seymour vertex $v$ lies in $S$.

Partition its deleted outneighbors as $A\dot\cup B$, where $B$ becomes
second neighbors and $A$ does not. Then $B\subseteq\bd T$. Put
\[
L=\bd(A\cup B)\setminus(S\cup T\cup\bd T).
\]
Every member of $L$ is a lost second neighbor of $v$: retained
intermediates lie in $S\cup T$, and neither can reach $L$ after deletion.
Thus the change in deficiency implies
$|L|-|B|<|A|+|B|$.
For $T'=T\cup A\cup B$,
\[
|\bd T'\setminus S|
 \le|\bd T\setminus S|+|L|-|B|
 <|T|+|A|+|B|=|T'|.
\]
The deleted outneighbor set is nonempty and disjoint from $T$.
This contradicts maximality. The empty-$T$ case uses
$|\bd T\setminus S|\le |T|=0$.
\end{proof}

In the proof, $L$ excludes $A\cup B$ through the definition of the
external boundary. This exclusion prevents original first neighbors
from being mistaken for lost second neighbors.

\begin{lemma}[Single-arc consequence; cf.\ \cite{BrantnerEtAl2009}]\label{lem:gap}
In an arc-minimal all-deficient graph, $g(x)\in\{1,2\}$ for every $x$.
\end{lemma}
\begin{proof}
Delete one arc $x\to y$. All other vertices retain their outneighbors,
and their second neighborhoods can only shrink, so the resulting
Seymour vertex must be $x$. Its outdegree decreases by one and its second
outdegree increases by at most one. Hence $g(x)\le2$.
\end{proof}

Lemma~\ref{lem:gap} is used for the refinements in
Section~\ref{sec:structure}; it is not needed for the main inequality.

\section{A target bundle and the local inequality}

Fix $x$ and define
\[
W=\Nin(x),\quad T=V(D)\setminus(\{x\}\cup W),\quad
P=\{v\ne x:x\notin\Nout(v)\cup\Ntwo(v)\}.
\]
The remaining part of $T$ is precisely $R=\Nback(x)$, so
$T=P\dot\cup R$. Every arc leaving $P$ ends in $T$.
Also
\begin{equation}\label{eq:T}
T=\Nout(x)\dot\cup\mathcal M(x).
\end{equation}

If $P=\varnothing$, then $\dback(x)=\dout(x)+\mu(x)$.
Since $x$ is deficient,
$\dtwo(x)-\dback(x)\le-\mu(x)-1\le\mu(x)-1$.

Assume henceforth $P\ne\varnothing$. Strong connectivity implies
$C:=\bd P\ne\varnothing$. Put
\[
c=|C|,\quad X=\bd^2P,\quad Z=R\setminus C,\quad
j=|Z|,\quad s=|\Nout(x)\cap Z|.
\]
We have $C\subseteq R$, $c=\dback(x)-j$, and
\begin{equation}\label{eq:X}
|X|\le c-1
\end{equation}
by Lemma~\ref{lem:boundary}.

Delete the $s$ arcs from $x$ to $Z$ and denote the resulting graph by
$D'$. When $s=0$, let $D'=D$.
The retained outneighbors of $x$ lie in $P\cup C$.
Every arc leaving $P\cup C$ has its head in $X$, while inside $P\cup C$
the vertices that are not retained outneighbors of $x$ are exactly its
original nonneighbors there. By \eqref{eq:T}, the set $Z$ contains
$j-s$ original nonneighbors. Consequently
\begin{equation}\label{eq:envelope}
d_{2,D'}^+(x)\le\mu(x)-(j-s)+|X|
 \le \mu(x)+\dback(x)-1-2j+s.
\end{equation}
A deleted outneighbor in $Z$ may become a second neighbor, but in that
case it already belongs to $X$. It must not be counted a second time.

If $s=0$, \eqref{eq:envelope} immediately proves \eqref{eq:main}.
If $s>0$, arc-minimality makes $x$ a Seymour vertex in $D'$:
all other vertices retain their outneighbors and cannot gain second
neighbors. Therefore
\[
\dout(x)-s\le d_{2,D'}^+(x)
 \le\mu(x)+\dback(x)-1-2j+s,
\]
and hence
\[
\dout(x)\le\mu(x)+\dback(x)-1-2(j-s).
\]
Using $\dtwo(x)\le\dout(x)-1$ proves the slightly stronger inequality
$\dtwo(x)-\dback(x)\le\mu(x)-2$ in this case.

This proves the local assertion of Theorem~\ref{thm:main}.
Finally, exact directed distance-two pairs are counted equally by
their sources and targets:
\[
\sum_x\dtwo(x)=\sum_x\dback(x).
\]
Summing \eqref{eq:main} gives $0\le2M-n$, completing the proof.


\section{Residuals, whole vertices, and near-equality}\label{sec:structure}

Define the nonnegative integer residual
\[
\eta(x)=\mu(x)-1+\dback(x)-\dtwo(x).
\]
Equivalently, with $p(x)=|P|$,
\[
\eta(x)=2\mu(x)+g(x)-1-p(x).
\]
We have the exact identity
\begin{equation}\label{eq:residual}
{\sum_x\eta(x)=2M-n.}
\end{equation}
For $P\ne\varnothing$, the preceding proof supplies the refinements
\[
s=0\Longrightarrow\eta(x)\ge2j,\qquad
s>0\Longrightarrow\eta(x)\ge g(x)+2(j-s).
\]

\begin{proposition}\label{prop:equality}
If $\eta(x)=0$, every nonneighbor of $x$ belongs to $\Ntwo(x)$.
Thus if $2M=n$, every missing pair is reachable in both directions by
paths of length two. More generally, at most $2M-n$ vertices have
positive residual.
\end{proposition}
\begin{proof}
For $P=\varnothing$, $\eta(x)=2\mu(x)+g(x)-1$, so zero residual forces
$\mu(x)=0$ and $g(x)=1$. Otherwise $s>0$ is impossible, and the
$s=0$ bound forces $j=0$.
Equality in \eqref{eq:envelope} and \eqref{eq:X} then gives
$|\Ntwo(x)|=\mu(x)+|X|$.
Here $\Ntwo(x)\subseteq\mathcal M(x)\dot\cup X$, since
$T=P\dot\cup C$. Therefore all nonneighbors are second outneighbors.
The final assertion follows from \eqref{eq:residual}.
\end{proof}

Call $x$ \emph{whole} when $\mu(x)=0$.

\begin{proposition}\label{prop:whole}
Every whole vertex of $D$ is reachable from every other vertex in at
most two steps. Consequently, if $w_2$ counts whole vertices with
$g(x)=2$, then $2M\ge n+w_2$.
\end{proposition}
\begin{proof}
The local inequality and Lemma~\ref{lem:gap} give
$p(x)\le g(x)-1\le1$.
Suppose $p(x)=1$. Then $g(x)=2$, $\eta(x)=0$, and the preceding proof
gives $P=\{u\}$, $j=s=0$ and
$\Nout(x)=\{u\}\dot\cup C$.
By definition $C=\Nout(u)$, so no vertex in $C$ points to $u$.
Thus the arc $x\to u$ has no alternative two-step path from $x$ to $u$.
Deleting it reduces the outdegree of $x$ by one but does not increase
its second outdegree. Vertex $x$ remains deficient, as do all other
vertices, contradicting arc-minimality.
Hence $p(x)=0$. A whole vertex now contributes $\eta(x)=g(x)-1$ to
\eqref{eq:residual}.
\end{proof}


\section{Sink components and the global missing-pair bound}\label{sec:sink}

The global bound extends beyond strongly connected graphs, with a
quantitative penalty for vertices outside a sink component. The local
inequality \eqref{eq:main} is still asserted only in the strongly
connected setting.

\begin{theorem}\label{thm:sink}
Let $D$ be any arc-minimal all-deficient oriented graph on $n$ vertices.
Let $S$ be a sink strongly connected component, put $k=|S|$, and put
$o=n-k$. Then
\begin{equation}\label{eq:sink}
M(D)\ge
\left\lceil\frac{k}{2}\right\rceil+
\left\lceil\frac{o(k-2)}2\right\rceil
\ge\left\lceil\frac n2\right\rceil+o.
\end{equation}
In particular, every arc-minimal counterexample satisfies
$M(D)\ge\lceil n/2\rceil$, without a connectivity assumption.
If equality holds, $D$ is strongly connected.
\end{theorem}
\begin{proof}
The graph $D[S]$ is all-deficient and arc-minimal. Indeed, deleting arcs
inside $S$ leaves every outside outneighborhood unchanged and can only
shrink its second neighborhood. A deletion preserving deficiency
throughout $S$ would therefore preserve it throughout $D$.
Theorem~\ref{thm:main} yields $M(D[S])\ge\lceil k/2\rceil$.

Write $O=V(D)\setminus S$, and let $\mu_S(v)$ count the missing pairs
from $v\in O$ to $S$. As $S$ is a sink, every adjacent pair from $O$ to
$S$ is directed into $S$. A two-step path entering $S$ cannot leave.
Consequently
\[
d_{1,D}^+(v)=d_{1,D[O]}^+(v)+k-\mu_S(v),\qquad
d_{2,D}^+(v)\le d_{2,D[O]}^+(v)+\mu_S(v),
\]
and hence
\[
g_D(v)\ge g_{D[O]}(v)+k-2\mu_S(v).
\]
For any oriented graph $H$,
$\sum_{v\in V(H)}g_H(v)\ge-2M(H)$: a present unordered pair supports at
most one distance-two direction, and a missing pair at most two.
Summing over $O$ and using $g_D(v)\le2$ gives
\[
2o\ge -2M(D[O])+ok-2M(O,S).
\]
Thus $M(D[O])+M(O,S)\ge\lceil o(k-2)/2\rceil$, proving the first bound.

Every all-deficient oriented graph has minimum outdegree at least two.
A sink is already a Seymour vertex; if $v$ has unique outneighbor $u$,
then every outneighbor of $u$ is an exact second neighbor of $v$, so
$v$ is again a Seymour vertex. The arc count on $D[S]$ therefore gives
$k\ge5$. Now
\[
\left\lceil\frac{k}{2}\right\rceil+
\left\lceil\frac{o(k-2)}2\right\rceil
\ge
\left\lceil\frac{k}{2}\right\rceil+
\left\lceil\frac{3o}{2}\right\rceil
\ge\left\lceil\frac n2\right\rceil+o.
\]
If $D$ is not strongly connected then $o\ge1$, proving the last assertion.
\end{proof}


\section{Uniform external out-neighbourhoods}
\begin{lemma}\label{lem:uniform}
Suppose $\varnothing\ne B\subsetneq V(D)$ and there is a set $C$ such
that $\Nout(v)\setminus B=C$ for every $v\in B$.
Then $D[B]$ has no arcs.
\end{lemma}
\begin{proof}
Strong connectivity gives $C=\bd B\ne\varnothing$.
Every vertex of $B$ points to every vertex of $C$, so no arc goes from
$C$ to $B$. Delete all arcs inside $B$. Each vertex of $B$ then has
first neighborhood $C$ and second neighborhood exactly $\bd^2B$.
Lemma~\ref{lem:boundary} makes it deficient. All other vertices retain
their outneighbors and cannot gain second neighbors. Thus a nonempty
internal arc set could be deleted without destroying the counterexample,
contrary to arc-minimality.
\end{proof}
